\section{What an executed interface result actually establishes}
\label{exec14:execution-meaning}

A proposal can contain an implemented local dispatcher, a proposed Java adapter,
a mocked external consumer and an exact mathematical reference at the same time.
The integration argument depends on keeping their evidence separate. The
scientific meaning of an output does not become stronger merely because it
passes through a real CLI or MCP server. A graph-syntax check remains a syntax
check; an unsupported policy request remains unsupported when its refusal is
correctly transported to a client.

The shared response therefore needs two independent descriptions. One records
how the result was obtained: design only, static inspection, project reference,
consumer double, actual local engine, actual transport, or an externally run
consumer. The other records the scientific disposition: computed within a
specified method, refused because declared premises conflict, unsupported,
unknown or unavailable. A generated dataset can be processed by an actual
engine. A mock can also return realistic numbers. Input provenance and execution
provenance must both be visible before a result is interpreted.

For an actual execution claim, the evidence binds the selected source and
schema, runtime and dependency versions, scientific context, input bytes,
route, request, response and terminal failure state. A run journal retains all
attempts and resource scope. An external package name or a source-code pin alone
is insufficient. Replays should be distinguished from fresh computations and
conflicting context should invalidate an idempotency key. Owner-run consumer
acceptance adds a named owner, actual consumer configuration and its acceptance
criteria; a local substitute cannot confer that acceptance.

The available local runtime implements nine narrow operations: capabilities,
graph syntax review, supplied-DAG back-door identification, supplied-moment
Gaussian residual inversion, deterministic quadratic sensitivity sets,
observation-clock diagnostics, reports and job status/cancellation. Its actual
local transport evidence includes eleven scientific cases agreeing across
library, CLI and an MCP SDK client/server exchange using protocol 2025-11-25.
The qualification applies to the recorded local stdio environment. It does not
establish modern-protocol or HTTP support, arbitrary discovery, a longitudinal
policy estimator or a live language-model backend. The local agent controller's
80 conformance checks involved zero live model calls
\citep{exec14_execution_evidence}.

A further servicing screen exercised 98 actual library and CLI invocations,
including restart/replay and changed-horizon conflicts. All 24 four-way semantic
comparisons agreed, and 964 independent rational membership checks reproduced
the returned deterministic sensitivity sets. The frozen suite nevertheless
retains a negative result: 94 of 98 expected outcomes matched. Four corrupt-input
cases returned the documented \texttt{artifact\_access} code where the authored
expectation named a different error. Each failed safely with no numerical
result or qualification. Diagnosing an expectation error does not rewrite the
original trial as a pass. The same trial shows that an actual engine can return
correctly scoped unsupported and refused results; it does not show that a
servicing estimator has been installed \citep{exec14_servicing_evidence}.

The distinction also changes cost interpretation. The servicing runner measured
35.4281 seconds summed across invocations and 36.5564 seconds of active runner
time. Parent and child peak resident memory were 45,056 and 104,448 KiB,
respectively; these are separate process maxima and cannot be added to claim a
concurrent memory peak. The timing covers this bounded local screen, not a
commercial engine, an external consumer or bank-scale deployment. A future
comparison should include setup, transfer, failures, review and recovery in the
same accounting boundary before proposing a service-level or economic benefit
\citep{exec14_servicing_evidence}.

\section{Servicing evidence, policy values and the proposed next experiment}
\label{exec14:servicing-evidence}

The servicing example has two different mathematical tasks. Exact enumeration
of a finite declared model can compare fixed policies. Estimating the value of
a policy from incomplete histories under hidden assignment information is a
harder problem. The available finite reference gives values $11/8$ for no
assistance, $2$ for assistance to everyone and $133/80$ under the stipulated
stress construction. These are generated model values, not fitted bank effects.
The library/CLI screen exercises supported graph, clock and sensitivity
operations; requests for the absent policy estimator return no policy estimate
\citep{exec14_servicing_evidence}.

The conservative precision calculation is an adverse result worth retaining.
At $n=1000$ the evaluation radius is approximately $8.242685524$, greater than
the complete single-policy range width of eight in the declared finite model.
The illustrative combined radius is approximately $9.615113642$. Such a bound
can be mathematically valid yet unhelpful for deciding between the fixed
policies. It is not evidence that collecting 1000 bank histories is a worthwhile
experiment. Better precision requires an explicitly justified sharper method
or a better information design; replacing the number with a narrow heuristic
would change the guarantee \citep{exec14_servicing_evidence}.

The missing-information examples are equally consequential. Two hidden-state
models can have identical observations but effects $.6$ and zero. Two
informative-dropout models can have identical observed histories but cohort
means $.25$ and $.75$. More observations of the same incomplete information do
not identify the missing distinction. A separate proposed continuous Gaussian
policy benchmark with quadratic reward also lies outside the bounded
finite-table theorem. It must retain its own tail and integrability analysis
rather than inherit the bounded guarantee by analogy
\citep{exec14_servicing_evidence}.

A planning scenario can still indicate what a prototype might be worth studying.
The fixed-policy contrast $2-11/8=5/8$, multiplied by 1000 customers and an
assumed USD100 per utility unit, gives USD62,500. Subtracting an assumed
additional full-workflow cost of USD10,000 gives USD52,500. Assistance costs
were already deducted inside the reward and are not subtracted again. Every
currency conversion, customer count and additional cost is stipulated. The
calculation is a scenario for sensitivity analysis, not measured bank ROI
\citep{exec14_servicing_evidence}.

The next experiment should first test a selected policy estimator against exact
references under controlled assignment, overlap and dropout mechanisms. It
should keep training and evaluation separate, retain every failed request, and
compare a fixed policy with the chosen trained-policy alternative under the same
reward and observation law. Bank data become relevant only if the prototype
isolates a consequential assumption or precision requirement that the public
and generated study cannot resolve. Until then, an honest negative precision
result and a typed refusal are useful proposal evidence.
